%
% 05_ableitung.tex -- Quaternionische Ableitung
%
% !TEX root = ../../buch.tex
% !TEX encoding = UTF-8
%
\section{Eine quaternionische Ableitung?\label{qa:section:ableitung}}
\kopfrechts{Quaternionische Ableitung}

Es ist naheliegend, die Differentialrechnung der komplexen Zahlen als Vorbild zu 
nehmen und zu fragen, ob sich dieser Aufbau auf $\mathbb{H}$ übertragen lässt. Die
Antwort ist eher enttäuschend, und der Grund dafür ist gerade jene Formel
$P' = qPq^{-1}$, die in Abschnitt~\ref{qa:section:drehungen} die Drehungen
möglich gemacht hat.

\subsection{Die komplexe Ableitung als Vorbild}
Eine Funktion $f\colon \mathbb{C} \to \mathbb{C}$ heisst in $z$ komplex
differenzierbar, wenn der Grenzwert
\begin{equation*}
    f'(z) = \lim_{h \to 0} \frac{f(z+h) - f(z)}{h}, \qquad h \in \mathbb{C}
    \setminus \{0\},
\end{equation*}
existiert. Entscheidend ist dabei, dass $h$ eine \emph{komplexe} Zahl ist und
sich der Null aus jeder beliebigen Richtung der Ebene nähern darf. Die
Forderung, dass der Grenzwert für alle diese Richtungen derselbe ist, ist eine
starke Bedingung. Sie hat zur Folge, dass eine einmal
komplex differenzierbare Funktion automatisch beliebig oft differenzierbar und
in eine Potenzreihe entwickelbar ist.

Zwei Dinge werden in dieser Definition stillschweigend benutzt: Erstens ist die
Division durch $h$ eindeutig, weil $\mathbb{C}$ kommutativ ist. Ob man
$h^{-1}$ von links oder von rechts multipliziert, spielt keine Rolle. Zweitens ist
der Ausdruck $h^{-1} z h = z$ für alle $h$, die Richtung von $h$ ist also für
die Lage von $z$ ohne Bedeutung. Beide Punkte gehen in $\mathbb{H}$ verloren.

\subsection{Zwei Differenzenquotienten}
Da die Multiplikation von Quaternionen nicht kommutativ ist, zerfällt der
Differenzenquotient in zwei verschiedene Ausdrücke. Man kann
\begin{equation*}
    \Delta_R(h) = \bigl(f(q+h) - f(q)\bigr)\,h^{-1}
    \qquad\text{oder}\qquad
    \Delta_L(h) = h^{-1}\,\bigl(f(q+h) - f(q)\bigr)
\end{equation*}
\index{Quaternionen!Differenzenquotienten}%
\index{Quaternionen!Ableitung}%
bilden, doch diese beiden Grössen stimmen im Allgemeinen nicht überein. Schon
die Schreibweise $\frac{f(q+h)-f(q)}{h}$ ist also nicht mehr eindeutig. Das
allein wäre noch kein grundsätzliches Hindernis, denn man könnte sich für eine der
beiden Varianten entscheiden. Im Folgenden wird mit $\Delta_R$ gearbeitet; für
$\Delta_L$ verläuft alles analog. Das eigentliche Problem liegt tiefer, wie das
folgende Beispiel zeigt.

\begin{beispiel}
    Untersucht wird die einfachste nichtlineare Funktion überhaupt,
    $f(q) = q^2$. Wegen $(q+h)^2 = q^2 + qh + hq + h^2$ ist die Differenz
    $f(q+h) - f(q) = qh + hq + h^2$, und der Differenzenquotient lautet
    \begin{equation*}
        \Delta_R(h) = (qh + hq + h^2)\,h^{-1}
        = q + h\,q\,h^{-1} + h .
    \end{equation*}
    Der erste und der letzte Term verhalten sich unproblematisch: $q$ ist
    konstant und $h \to 0$. Der mittlere Term $h q h^{-1}$ ist jedoch genau die
    Drehformel aus Abschnitt~\ref{qa:section:drehungen:drehformel}. Schreibt man
    $h = t \cdot \hat{h}$ mit einer reellen Zahl $t > 0$ und einer 
    Einheitsquaternion
    $\hat{h}$, so kürzt sich der reelle Faktor heraus,
    \begin{equation*}
        h\,q\,h^{-1} 
        = t\,\hat{h}\,q\,\hat{h}^{-1}\,t^{-1} 
        = \hat{h}\,q\,\hat{h}^{-1},
    \end{equation*}
    und übrig bleibt die Drehung von $q$ um die Achse $\operatorname{Im}(\hat{h})$.
    Dieser Term wird also nicht klein, wenn $h$ klein wird, sondern hängt nur von 
    der Richtung ab, in der $h$ gegen null geht.
\end{beispiel}



Um das Beispiel konkret zu machen, sei $q = k$, also der Punkt $(0,0,1)$. Nähert 
man sich der Null entlang
der reellen Achse, $h = t \in \mathbb{R}$, so ist $hqh^{-1} = q$ und
\begin{equation*}
    \Delta_R(t) 
    = k + k + t \;\rightarrow\; 2k .
\end{equation*}
Nähert man sich dagegen entlang der $i$-Achse, $h = ti$, so ist
$i k i^{-1} = -k$, denn nach Abschnitt~\ref{qa:section:drehungen} bewirkt $i$
eine Drehung um $180^\circ$ um die $x$-Achse, und diese führt $(0,0,1)$ in
$(0,0,-1)$ über. Damit ist
\begin{equation*}
    \Delta_R(ti) 
    = k - k + ti \;\rightarrow\; 0 .
\end{equation*}
Die beiden Grenzwerte $2k$ und $0$ sind verschieden, der Grenzwert existiert
also nicht. Durchläuft $u$ alle Einheitsquaternionen, so durchläuft
$\hat{h}\,q\,\hat{h}^{-1}$ die gesamte Kugel vom Radius $|\operatorname{Im}(q)|$ 
um den Realteil von $q$; die möglichen Richtungsgrenzwerte bilden eine ganze
Kugelschale und nicht einen einzelnen Punkt. Sie fallen genau dann zusammen,
wenn $\operatorname{Im}(q) = 0$ ist, wenn also jede Drehung $q$ festhält. Die
Funktion $q^2$ ist somit nur in den reellen Punkten differenzierbar, und dort
ergibt sich mit $\Delta_R \to 2q$ erwartungsgemäss die reelle Ableitung.

\subsection{Fazit}
Quaternionen sind ein hervorragendes
\emph{algebraisches} Werkzeug, um Drehungen zu beschreiben und zu verketten. Für 
Differenzialrechnungen sind sie jedoch nicht geeignet. Wo man in einer Anwendung 
Ableitungen benötigt, etwa bei
der Geschwindigkeitsbestimmung von Drehungen, arbeitet man deshalb
nicht mit einer quaternionischen Ableitung, sondern differenziert die vier
reellen Komponenten einzeln.
